% Generated by build_parameter_table.py; do not hand-edit.
\clearpage\begingroup
\captionof{table}{The six anchoring/fitting studies in the current joint objective. All rows are in-sample, outcome-informed calibration, not independent validation. Residual counts are numerical targets, not independent animals. All nine fitted coefficients/proxies were optimized jointly; no one-to-one study-to-parameter identifiability is claimed.}\label{tab:fitting-studies}
\centering\fontsize{8}{10}\selectfont
\setlength{\tabcolsep}{4pt}\renewcommand{\arraystretch}{1.12}
\begin{tabular}{@{}>{\raggedright\arraybackslash}p{2.7cm}>{\raggedright\arraybackslash}p{3.7cm}>{\raggedright\arraybackslash}p{4.3cm}r>{\raggedright\arraybackslash}p{5.6cm}@{}}
\toprule \textbf{Study / species} & \textbf{Protocol represented} & \textbf{Fitted observations / residual scaling} & \textbf{Count} & \textbf{Transfer / observation assumptions} \\ \midrule
Pike and Bosher (1980)~\citep{PikeBosher1980}\newline Guinea pig & Furosemide, 40/60/80 mg/kg; implemented infusion 0.542 min. & Baseline-relative EP curves; errors divided by \(10\,\mathrm{mV}\sqrt{n_c}\). & 37 & Manual curve anchors; effective dose proxy, not measured drug exposure. \\[4pt]
Asakuma and Snow (1980)~\citep{AsakumaSnow1980}\newline Guinea pig & 60 mg/kg furosemide; implemented infusion 0.5 min. & EP reduction and resistance ratio at modeled EP nadir; scales 10 mV and 0.2. & 2 & Difference/ratio of rounded group means; 0--120-min model nadir is an assumed observation operator. \\[4pt]
Rybak and Morizono (1982)~\citep{RybakMorizono1982}\newline Chinchilla & 25/50/100 mg/kg furosemide; assumed infusion 0.5 min. & EP and endolymph-K reductions; each error scaled by its reported SE. & 6 & EP and K minima selected separately over 0--120 min; group-mean minima need not equal minima of a representative trajectory. \\[4pt]
Wang et al. (1992)~\citep{Wang1992}\newline Guinea pig & 105/115/125 dB SPL noise, 20 min, then quiet; anesthetized recordings. & Baseline-relative EP curves; errors divided by \(3\,\mathrm{mV}\sqrt{n_c}\). & 45 & Broadband level-to-drive map is empirical; digitized group curves have unknown covariance. \\[4pt]
Wang et al. (1990)~\citep{Wang1990}\newline Guinea pig & 125 dB SPL for 20/40/60/80 min in awake restrained animals. & Exposed-minus-control cohort EP differences; errors divided by \(5\,\mathrm{mV}\sqrt{4}\). & 4 & Protocol mismatch: observations include anoxia, which is not modeled. A fixed 45-min post-noise delay is assumed; source text/figure counts differ. \\[4pt]
Salt and Konishi (1979)~\citep{SaltKonishi1979}\newline Guinea pig & Table III: three 20-min exposures at 95/105/115 dBA, with 5-min quiet intervals. Fig. 3: 115 dBA at 0--20, 30--50 and 60--80 min. & Table III local EP changes: \(3\,\mathrm{mV}\sqrt{3}\) scale. Fig. 3 paired EP/K curves: \(3\,\mathrm{mV}\sqrt{n_c}\) and \(3\,\mathrm{mM}\sqrt{n_c}\). & 9 + 75 & Same-number dBA-to-model-SPL mapping; Table III 1--3-min summary window is assumed, n counts observations, not animals. Fig. 3 is one animal; Fig. 5 excluded. \\[4pt]
\midrule\multicolumn{3}{l}{\textbf{Total: six studies, seven residual components}} & 178 & Salt contributes two components, not two independent studies. \\
\bottomrule\end{tabular}
\par\smallskip\raggedright\footnotesize \(n_c\) is the number of retained digitized anchors in the corresponding curve, not an animal sample size. Deterministic scales balance endpoints/curves; even SE-scaled terms do not define a joint likelihood when covariance and discrepancy are unknown. Source SD/SE and available sample sizes are retained in Figure~\ref{fig:tier1} and the calibration data. Salt Figure 5 is excluded. Zhang (2021) and Hou (2020) are retrospective conditional benchmarks only, not additional fitting studies. Infusion lengths and timing windows above describe the implemented observation operators, with assumptions explicitly marked.
\par\endgroup
